%! TeX program = lualatex
\documentclass[twocolumn]{article}

\usepackage{xfrac}
\usepackage[mathletters]{ucs}
\usepackage[margin=0.1in]{geometry}
\usepackage{amsmath,amssymb,amsfonts}
\author{}
\setlength{\topsep}{0em}
\raggedbottom
\raggedright

\def\epsilon{\varepsilon}
\NewCommandCopy{\oldbf}{\textbf}
\def\colo{white}
\renewcommand{\textbf}[1]{\textcolor{\colo}{\oldbf{#1}}}
% supports augmented matrix
\newenvironment{mat}[1]{\left[\begin{array}{#1}}{\end{array}\right]}
\newenvironment{vmat}[1]{\left|\begin{array}{#1}}{\end{array}\right|}

% no margin align
\makeatletter
\renewenvironment{align*}{\ifvmode\else\hfil\null\linebreak\fi
  \hspace*{-\leftmargin}\minipage\textwidth
  \setlength{\abovedisplayskip}{0pt}
  \setlength{\abovedisplayshortskip}{\abovedisplayskip}
  \start@align\@ne\st@rredtrue\m@ne}
{\endalign\endminipage\linebreak}
\makeatother

\usepackage{tikz}
\usetikzlibrary{angles,quotes}

\usepackage{listings}
\lstset{
	basicstyle={\ttfamily},
	commentstyle=\color{gray},
	keywordstyle=\color{pink},
	columns=flexible,
}

\usepackage[no-math]{fontspec}
\setsansfont{SourceSansPro}
\setmonofont{Iosevka Custom}
\renewcommand{\familydefault}{\sfdefault}

\usepackage[UKenglish]{datetime}
\renewcommand{\dateseparator}{-}
\yyyymmdddate
\date{\today}

\usepackage{xcolor}
\definecolor{black}{HTML}{000000}
\definecolor{darkest}{HTML}{060606}
\definecolor{darker}{HTML}{131313}
\definecolor{dark}{HTML}{313131}
\definecolor{gray}{HTML}{787878}
\definecolor{light}{HTML}{bfbfbf}
\definecolor{lighter}{HTML}{dddddd}
\definecolor{lightest}{HTML}{eaeaea}
\definecolor{white}{HTML}{f1f1f1}
\definecolor{red}{HTML}{c53943}
\definecolor{green}{HTML}{39c575}
\definecolor{orange}{HTML}{ff9966}
\definecolor{blue}{HTML}{3989c5}
\definecolor{indigo}{HTML}{3943c5}
\definecolor{leek}{HTML}{39c5bb}
\definecolor{pink}{HTML}{ff6680}
\definecolor{lime}{HTML}{66ff99}
\definecolor{yellow}{HTML}{ffe666}
\definecolor{sky}{HTML}{66ccff}
\definecolor{purple}{HTML}{a68cd9}
\definecolor{cyan}{HTML}{66ffe6}
\colorlet{fg}{white}
\colorlet{bg}{black}
\pagecolor{bg}
\color{fg}

\usepackage{hyperref}
\hypersetup{colorlinks=true,linktoc=all,linkcolor=yellow}

\usepackage{mdframed}
\mdfdefinestyle{default}{
	topline=false,
	bottomline=false,
	rightline=false,
	linewidth=3pt,
	fontcolor=fg,
	backgroundcolor=bg,
	% nobreak=true,
}
% \newenvironment{framed}[3]{
% 	\def\colo{#1}
% 	\begin{minipage}{\linewidth}
% 		\begin{mdframed}[style=default,linecolor=#1,backgroundcolor=#1!20!bg]
% 			\textcolor{#1}{\textbf{#2.}} #3
% 		\end{mdframed}
% 		\begin{mdframed}[style=default,linecolor=#1,backgroundcolor=#1!10!bg]
% }{\end{mdframed}\end{minipage}\def\colo{white}}

\newenvironment{framed}[3]{
	\def\colo{#1}
		\begin{mdframed}[style=default,linecolor=#1,backgroundcolor=#1!20!bg]
			\textcolor{#1}{\textbf{#2.}} #3
		\end{mdframed}
		\begin{mdframed}[style=default,linecolor=#1,backgroundcolor=#1!10!bg]
}{\end{mdframed}\def\colo{white}}

\newenvironment{theorem}[1][]{\begin{framed}{leek}{Theorem}{#1}}{\end{framed}}
\newenvironment{definition}[1][]{\begin{framed}{sky}{Definition}{#1}}{\end{framed}}
\newenvironment{algorithm}[1][]{\begin{framed}{pink}{Algorithm}{#1}}{\end{framed}}
\newenvironment{example}[1][]{\begin{framed}{orange}{Example}{#1}}{\end{framed}}

\newenvironment{leek}[1][]{\begin{tframed}{leek}{#1}{}}{\end{tframed}}


\lstset{
  language=C++,
  morekeywords={in},
}
\title{A22 | Linear Algebra}
\begin{document}

\maketitle
\tableofcontents
\section{Complex Numbers}
    \subsection{Cartesian}
		\begin{definition}
			$z=a+bi\quad w=c+di$ for $a,b,c,d\in\mathbb{R}$
			\begin{itemize}
				\item $a=Re(z)=r\cos\theta$ is the \textbf{real} part of $z$
				\item $b=Im(z)=r\sin\theta$ is the \textbf{imaginary} part of $z$
				\item $z=w$ iff $a=c\land b=d$
			\end{itemize}
		\end{definition}
	\subsection{Polar}
	    \begin{definition}
	    	$z=re^{i\theta}\quad w=se^{i\varphi}$ for $r,s,\theta,\varphi\in\mathbb{R}$
			\begin{itemize}
				\item $r=|z|=\sqrt{a^2+b^2}$ is the \textbf{modulus} of $z$
				\item $\theta=\arg z=\mathrm{atan}(\sfrac ba)$ is \textbf{an argument} of $z$\\
					$\theta\in(-\pi,\pi]$ is \textbf{the principle argument} of $z$\\
					$\mathrm{atan}(\sfrac{-b}a)$ and $\mathrm{atan}(\sfrac b{-a})$ are not the same for $a,b>0$
				\item $z=re^{i(\theta+2\pi k)}$ is the \textbf{general solution}
			\end{itemize}
	    \end{definition}
	\subsection{Add}
		\begin{definition}[$z+w=(a+c)+(c+d)i$]
			 \begin{itemize}
				\item commutativity, associativity
				\item add-identity: $0$
				\item add-inverse: $-a-bi=-re^{i\theta}$
			 \end{itemize}
		\end{definition}
	\subsection{Mul}
		\begin{definition}[$zw=rse^{i(\theta+\varphi)}$]
			\begin{itemize}
				\item commutativity, associativity, distributivity
				\item mul-identity: $1$
				\item mul-inverse: $\displaystyle\frac1z=\frac{\overline z}{|z|^2}$
			\end{itemize}
		\end{definition}
	\subsection{Conjugate}    
		\begin{definition}[$\overline z=a-bi$]
			 \begin{itemize}
				\item $\overline{z\pm w}=\overline z\pm\overline w$ 
				\item $\overline{zw}=\bar z\bar w$
				\item $\overline{(\bar z)}=z$
				\item $\overline z=z$ iff $z\in\mathbb{R}$
				\item $z\overline z=|z|^2=r^2=a^2+b^2$
			 \end{itemize}
		\end{definition}
	\subsection{Modulus}
		\begin{definition}
			$|z|=\sqrt{a^2+b^2}=r$
			\begin{itemize}
				\item $|z|^2=z\overline z=a^2+b^2$
				\item $|zw|=|z||w|$
				\item $|z+w|\le|z|+|w|$ ($\Delta$ inequality)
				\item $|z-w|=\sqrt{(a-c)^2+(b-d)^2}$ (distance)
			\end{itemize}
		\end{definition}
		\begin{example}[Prove $|z|\le|a|+|b|\le\sqrt2|z|$ for $z=a+bi\in\mathbb{C}$]
			\begin{align*}
				|z|&\le|a|+|b|\\
				|z|^2&\le |a|^2+|b|^2+2|a||b|\\
				a^2+b^2&\le a^2+b^2+2|ab|\\
				0&\le|ab|
			\end{align*}
			\begin{align*}
				|a|+|b|&\le\sqrt2|z|=\sqrt{2a^2+2b^2}\\
				|a|^2+|b|^2+2|a||b|&\le2a^2+2b^2\\
				0&\le |a|^2+|b|^2-2|a||b|\\
				0&\le(|a|-|b|)^2
			\end{align*}
		\end{example}
		\begin{example}[Prove $zw\in\mathbb{R}\land z\ne0\implies\exists a\in\mathbb{R}$ st $w=a\overline z$]
			assume $zw\in\mathbb{R}$ and $z\ne0$, let $z=re^{i\theta}\enspace w=se^{i\varphi}$
			\begin{align*}
				\therefore zw&=rse^{i(\theta+\varphi)}=rs\tag*{$zw\in\mathbb{R}\implies i=0$}\\
				w&=\frac{rs}z=\frac{rs\overline z}{|z|^2}\tag*{$z\ne0\quad\frac1z=\frac{\overline z}{|z|^2}$}\\
				 &=\frac{rs\overline z}{r^2}=\frac sr\overline z=az\quad\blacksquare\tag*{$z=ne0\implies r\ne0$}\\
			\end{align*}
		\end{example}
	\subsection{Euler $\times$ De Moivre's}
		\begin{theorem}[Euler's Formula]
			\begin{align*}
				e^{i\theta}&=\cos\theta+i\sin\theta\\
				\implies e^{in\theta}&=\cos(n\theta)+i\sin(n\theta)
			\end{align*}
		\end{theorem}
		\begin{theorem}[De Moivre's]
			\begin{align*}
				(e^{i\theta})^n&=e^{in\theta}\\
				\implies (\cos\theta+i\sin\theta)^n&=\cos(n\theta)+i\sin(n\theta)\tag*{from Euler's}
			\end{align*}
		\end{theorem}
		\begin{example}[$\cos2\theta=\cos^2\theta-\sin^2\theta\quad\sin3\theta=3\cos^2\theta\sin\theta-\sin^3\theta$]
			\begin{align*}
				\cos2\theta&=Re(\cos2\theta+i\sin2\theta)\\
				&=Re((\cos\theta+i\sin\theta)^2)\tag*{De Moivre}\\
				&=Re(\cos^2\theta-\sin^2\theta+2i\sin\theta\cos\theta)\\
				&=\cos^2\theta-\sin^2\theta\\
				\sin3\theta&=Im(\cos3\theta+i\sin3\theta)\\
			    &=Im((\cos\theta+i\sin\theta)^3)\tag*{De Moivre}\\
				&=3\cos^2\theta\sin\theta-\sin^3\theta\\
			\end{align*}
		\end{example}
	\subsection{Roots}
		\begin{theorem}[Roots]
			\begin{itemize}
			   \item if $z,q\in\mathbb{C},n\in\mathbb{Z}^+$ and $z^n=q$, then $z$ is \textbf{nth root} of $q$
			   \item $z^n=r^ne^{ni(\theta+2k\pi)}$ has $n$ solutions for $0\le k\le n$
			\end{itemize}
		\end{theorem}
		\begin{example}[Solve $z^6=-64$]
			let $z=re^{i\theta}$
			\begin{align*}
				z^6&=(re^{i\theta})^6=r^6e^{6i\theta}=64e^{i\pi}\\
				\therefore r^6&=64\implies r=2\\
				6\theta&=\pi+2\pi k\quad 0\le k\le6\\
				\theta&=\frac{\pi+2\pi k}6\\
				\therefore z&=2e^{i(\sfrac{\pi}6)},2e^{i(\sfrac{3\pi}6)},2e^{i(\sfrac{5\pi}6)},2e^{i(\sfrac{7\pi}6)},2e^{i(\sfrac{9\pi}6)},2e^{i(\sfrac{11\pi}6)}
			\end{align*}
		\end{example}
\section{Systems of Linear Equations}
	\begin{definition}[System of $m$ Equations with $n$ Variables]
		$$A\vec x=\vec b$$
		where 
		\begin{itemize}
			\item $A\in\mathbb{R}^{m\times n}$ is the \textbf{coefficient} matrix
			\item $\vec x\in\mathbb{R}^n$ are the vairables
			\item $\vec b\in\mathbb{R}^m$ is the \textbf{constant vector}
			\item $\mathrm{rank}\begin{mat}{cc}A&b\end{mat}=\mathrm{rank}(A)\implies$ system is consistent
			\item $A$ is invertible $\implies$ system has unique solution $\forall\vec b$
		\end{itemize}
	\end{definition}
    \subsection{Gaussian Elimination}
	    \subsubsection{Elementary Row Operations}
		    \begin{definition}
				\begin{enumerate}
					\item interchange two rows\quad\texttt{swap(ai, aj)}
					\item multily a row by a nonzero number\quad\texttt{ai *= n}
					\item add a multiple of one row to another\quad\texttt{ai += n * aj}
				\end{enumerate}
		    \end{definition}
        \subsubsection{Row-Echelon Form}
            \begin{definition}[REF]
                \begin{itemize}
                    \item pivot $\equiv$ leading 1
                    \item each row has pivot $\oplus$ is all 0
                    \item each pivot is to the right of all pivots in the rows above it
					\item all zero-rows are at the bottom
                \end{itemize}
            \end{definition}
            \begin{example}[Matrices in REF]
                (I can't draw the stairs)
                $$
                \begin{mat}{ccc}1&*&*\\0&1&*\end{mat}
                \begin{mat}{cccc}0&1&*&*\\0&0&1&*\\0&0&0&0\end{mat}
                \begin{mat}{cccc}1&*&*&*\\0&1&*&*\\0&0&0&1\end{mat}
                \begin{mat}{ccc}1&*&*\\0&1&*\\0&0&1\end{mat}\\
                $$
            \end{example}
            \begin{definition}[Reduced Row-Echelon Form (RREF)]
                \begin{itemize}
                    \item is in REF
                    \item all entries above a pivot is 0
                \end{itemize}
            \end{definition}
        \subsubsection{Gaussian Elimination}
            \begin{algorithm}
                \begin{lstlisting}
for c, col in a.cols.enumerate()
    for r, row in a.rows.enumerate()
        if col[r] == 0 
            continue
        r /= a[r,c] // make pivot
        move row to top
    for r, row in a.rows.enumerate().skip(1)
        row -= a[r,c] * a.top // make zeros
                \end{lstlisting}
				\begin{itemize}
					\item if has $\begin{mat}{ccc|c}
						0 & \cdots & 0 & 1\\
					\end{mat}$, then is inconsistent
				\end{itemize}
            \end{algorithm}
            \begin{example}[No Solution]
                \begin{align*}
                    \begin{mat}{ccc|c}
                        \textcolor{red}3 & \textcolor{red}1 & \textcolor{red}{-4} & \textcolor{red}{-1}\\
                        \textcolor{lime}1 & \textcolor{lime}0 & \textcolor{lime}{10} & \textcolor{lime}5\\
                        \textcolor{sky}4 & \textcolor{sky}1 & \textcolor{sky}6 & \textcolor{sky}1\\
                    \end{mat}&\sim\begin{mat}{ccc|c}
                        \textcolor{lime}1 & \textcolor{lime}0 & \textcolor{lime}{10} & \textcolor{lime}5\\
                        \textcolor{red}0 & \textcolor{red}{-1} & \textcolor{red}{34} & \textcolor{red}{16}\\
                        \textcolor{sky}0 & \textcolor{sky}{-1} & \textcolor{sky}{34} & \textcolor{sky}{19}\\
                    \end{mat}\\
                    &\sim\begin{mat}{ccc|c}
                        \textcolor{lime}1 & \textcolor{lime}0 & \textcolor{lime}{10} & \textcolor{lime}5\\
                        \textcolor{red}0 & \textcolor{red}{-1} & \textcolor{red}{34} & \textcolor{red}{16}\\
                        \textcolor{sky}0 & \textcolor{sky}0 & \textcolor{sky}0 & \textcolor{sky}{-3}\\
                    \end{mat}
                \end{align*}
                $0=-3\implies$has no solution, is inconsistent
            \end{example}
            \begin{example}[Has Parameters (General Solution)]
                \begin{align*}
                    \begin{mat}{cccc|c}
                        1 & -2 & -1 & 3 & 1\\
                        \textcolor{lime}2 & \textcolor{lime}{-4} & \textcolor{lime}1 & \textcolor{lime}0 & \textcolor{lime}5\\
                        \textcolor{sky}1 & \textcolor{sky}{-2} & \textcolor{sky}2 & \textcolor{sky}{-3} & \textcolor{sky}4\\
                    \end{mat}&\sim\begin{mat}{cccc|c}
                        1 & -2 & -1 & 3 & 1\\
                        \textcolor{lime}0 & \textcolor{lime}0 & \textcolor{lime}{-3} & \textcolor{lime}{-6} & \textcolor{lime}3\\
                        \textcolor{sky}0 & \textcolor{sky}0 & \textcolor{sky}3 & \textcolor{sky}{-6} & \textcolor{sky}3\\
                    \end{mat}\\
                    &\sim\begin{mat}{cccc|c}
                        1 & -2 & -1 & 3 & 1\\
                        \textcolor{lime}0 & \textcolor{lime}0 & \textcolor{lime}1 & \textcolor{lime}{-2} & \textcolor{lime}1\\
                        \textcolor{sky}0 & \textcolor{sky}0 & \textcolor{sky}0 & \textcolor{sky}0 & \textcolor{sky}0\\
                    \end{mat}\\
                    &\sim\begin{mat}{cccc|c}
                        \textcolor{red}1 & \textcolor{red}{-2} & \textcolor{red}0 & \textcolor{red}1 & \textcolor{red}2\\
                        0 & 0 & 1 & -2 & 1\\
                        0 & 0 & 0 & 0 & 0\\
                    \end{mat}
                \end{align*}
                no pivots for $x_2$ and $x_4\quad\implies x_2=s\quad x_4=t$
                \begin{align*}
                    x_1&=2+2s-t\\
                    x_2&=s\\
                    x_3&=1+2t\\
                    x_4&=t
                \end{align*}
                for $s,t\in\mathbb{R}$
            \end{example}
            \begin{example}[Has Conditions]
                \begin{align*}
                    \begin{mat}{ccc|c}
                        1 & 3 & 1 & a\\
                        -1 & -2 & 1 & b\\
                        3 & 7 & -1 & c\\
                    \end{mat}&\sim\begin{mat}{ccc|c}
                        1 & 3 & 1 & a\\
                        0 & 1 & 2 & a+b\\
                        0 & -2 & -4 & c-3a\\
                    \end{mat}\\
                    &\sim\begin{mat}{ccc|c}
                        1 & 0 & -5 & -2a-3b\\
                        0 & 1 & 2 & a+b\\
                        0 & 0 & 0 & c-a+2b\\
                    \end{mat}
                \end{align*}
                $0=c-a+2b\implies c=a-2b$\\
                if $c=a-2b$, then
                \begin{align*}
                    x_1&=5t-2a-3b\\
                    x_2&=a+b-2t\\
                    x_3&=t
                \end{align*}
				for $t\in\mathbb{R}$
            \end{example}
        \subsubsection{Rank of a Matrix}\label{sec:rank}
            \begin{definition}
                number of pivots after being Gaussian Eliminated
            \end{definition}\,\\
            \begin{theorem}
                if a $m\times n$ ($m$ rows $\times n$ cols) matrix is \textbf{consistent}, and it has rank $r$, then
                \begin{itemize}
                    \item the solution has $n-r$ parameters
                    \item if $r<n$, has $\infty$ solutions
                    \item if $r=n$ has a unique solution
                \end{itemize}
            \end{theorem}
    \subsection{Homogeneous Equations}
        \begin{definition}
			$$A\vec x=\vec 0$$
            \begin{itemize}
                \item $x_i=0$ is the \textbf{trivial solution}
                \item if parameter $\exists$, the system has a general \textbf{non trivial solution} ($\therefore$ inf solutions)
            \end{itemize}
        \end{definition}
        \begin{theorem}[$A\vec x=\vec0$ for $A\in\mathbb{R}^{m\times n}$]
            if $m<n$, then it has a nontrival solution\\
            \textbf{Proof.} suppose a RREF augmented matrix $A$ has $m$ rows and $n$ cols, $n>m$
            \begin{align*}
                \text{suppose }&A\text{ has }r\text{ leading vars }(0\le r\le m)\\
                &\therefore A\text{ has }n-r\text{ non leading vars (parameters)}\\
                &\because m<n\\
                &\therefore 0\le r<n\\
                &\therefore n-r>0\\
                &\therefore A\text{ has parameters}\quad\blacksquare
            \end{align*}
        \end{theorem}
        \subsubsection{Solution as Linear Combination}
            \begin{definition}
				suppose $\vec x,\vec y,\vec z\in\mathbb{R}^n$, $\vec z$ is a \textbf{linear combination} of $x,y$ if
				$$\vec z=s\vec x+t\vec y$$
				for $s,t\in\mathbb{R}$
				\begin{itemize}
					\item $\vec 0$ is a linear combination of all vectors
				\end{itemize}
            \end{definition}
            \begin{example}[Find general solution as linear comb of basic solutions]
                $$A=\begin{mat}{cccc}
                    1 & -2 & 3 & -2\\
                    -3 & 6 & 1 & 0\\
                    -2 & 4 & 4 & -2\\
                \end{mat}$$
                \begin{enumerate}
                    \item convert to RREF augmented
                        $$\begin{mat}{cccc|c}
                            1 & -2 & 3 & -2 & 0\\
                            -3 & 6 & 1 & 0 & 0\\
                            -2 & 4 & 4 & -2 & 0\\
                        \end{mat}\sim\begin{mat}{cccc|c}
                            1 & -2 & 0 & -\sfrac15 & 0\\
                            0 & 0 & 1 & -\sfrac35 & 0\\
                            0 & 0 & 0 & 0 & 0\\
                        \end{mat}$$
                    \item general solution
                        $$x_1=2s+\sfrac15t\quad x_2=s\quad x_3=\sfrac35t\quad x_4=t$$
                        $$\vec x
                        =\begin{mat}{c}
                            2s+\sfrac15t\\ s\\ \sfrac35t\\ t\\
                        \end{mat}
                        =s\begin{mat}{c}
                            2\\ 1\\ 0\\ 0\\
                        \end{mat}+t\begin{mat}{c}
                            \sfrac15\\ 0\\ \sfrac35\\ 1\\
                        \end{mat}=s\vec x_1+t\vec x_2$$
					\item $\vec x_1$ and $\vec x_2$ are basic solutions
                \end{enumerate}
            \end{example}
		\subsubsection{Solution of Any System as $\vec x_p+\vec x_0$}
			\begin{algorithm}[Solve $A\vec x=\vec b$ as $\vec x_p$ and $\vec x_0$]
				\begin{itemize}
					\item convert to RREF
					\item $\vec x_p$ is solution when all parameters are $0$
					\item $\vec x_0$ is the general solution when $\vec b=\vec0$
				\end{itemize}
			\end{algorithm}
			\begin{example}
					$$\begin{mat}{cccc|c}
						2 & 1 & -1 & -1 & -1\\
						3 & 1 & 1 & -2 & -2\\
						-1 & -1 & 2 & 1 & 2\\
						-2 & -1 & 0 & 2 & 3\\
					\end{mat}\sim\begin{mat}{cccc|c}
						1 & 0 & 0 & \textcolor{sky}1 & \textcolor{leek}3\\
						0 & 1 & 0 & \textcolor{sky}{-4} & \textcolor{leek}{-9}\\
						0 & 0 & 1 & \textcolor{sky}{-1} & \textcolor{leek}{-2}\\
						0 & 0 & 0 & \textcolor{sky}0 & \textcolor{leek}0\\
					\end{mat}$$
					$$\vec x_p=\begin{mat}{c} \textcolor{leek}3\\\textcolor{leek}{-9}\\\textcolor{leek}{-2}\\\textcolor{leek}0\end{mat}\quad\vec x_0=\begin{mat}{c} \textcolor{sky}{-1}\\ \textcolor{sky}4\\ \textcolor{sky}1\\ \textcolor{sky}0\\ \end{mat}x_4$$
					$$\vec x=\begin{mat}{c} 3\\ -9\\ -2\\ 0\\ \end{mat}+\begin{mat}{c} -1\\ 4\\ 1\\ 0\\ \end{mat}x_4$$
			\end{example}



\section{Matrix Algebra}
    \begin{definition}
        \begin{itemize}
            \item $m\times n$ matrix means $m$ rows and $n$ cols
            \item $a_{ij}$ means at row $i$ and col $j$
        \end{itemize}
    \end{definition}
	\subsection{Properties}
        \subsubsection{Identity}	
			\begin{definition}[Identity Matrix]
				$$I_1=\begin{mat}{c}1\end{mat}\quad 
				I_2=\begin{mat}{cc}
					1 & 0\\
					0 & 1\\
				\end{mat}\quad
				I_n=[e_{ij}]\text{ where } e_{ij}=\begin{cases}
					1&i=j\\
					0&i\ne j\\
				\end{cases}$$
				\begin{itemize}
					\item $IA = AI = A$
				\end{itemize}
			\end{definition}
			\begin{definition}[Standard Basis of $R^n$]
				$I_n=\begin{mat}{cccc} e_1 & e_2 & \cdots & e_n\\ \end{mat}\quad A=\begin{mat}{cccc} a_1 & a_2 & \cdots & a_n\end{mat}$
				\begin{itemize}
					\item $Ae_i=a_i$ 
					\item useful for proofs in \hyperref[sec:standard-basis]{Mat-Vec Mul}
				\end{itemize}
			\end{definition}
		\subsubsection{(Skew) Symmetry}
			\begin{definition}
				\begin{itemize}
					\item $A$ is \textbf{symmetric} iif $A^T=A$
					\item $A$ is \textbf{skew symmetric} iif $A^T=-A$
				\end{itemize}
			\end{definition}
			\begin{example}[Prove $\forall A\in\mathbb{R}^{n\times n}\exists$ symmetric $S$ and skew symmetric $W$ st $A=S+W$ and are uniquely determined by $A$]
				let $A=S+W$ where $S^T=S$ and $W^T=-W$
				\begin{align*}
					\therefore A^T&=S^T+W^T=S-W\\
					A+A^T&=2S\implies S=\sfrac12(A+A^T)\\
					A-A^T&=2W\implies W=\sfrac12(A-A^T)
				\end{align*}
			\end{example}
		\subsubsection{Idempotent}
		    \begin{definition}
		    	for $A\in\mathbb{R}^{n\times n}$, $A$ is \textbf{idempotent} iif
				$$A^2=A$$
		    \end{definition}
			\begin{example}[Prove $\forall A\in\mathbb{R}^{n\times n}$\\$P$ is idempotent $\implies Q=P+AP-PAP$ is idempotent]
				assume $P=P^2$
				\begin{align*}
					Q&=P+AP-PAP\\
					\because\textcolor{\colo}PQ&=\textcolor{\colo}{P^2}+\textcolor{\colo}PAP-\textcolor{\colo}{P^2}AP=P+PAP-PAP=P\\
					\therefore Q^2&=P\textcolor{\colo}Q+AP\textcolor{\colo}Q-PAP\textcolor{\colo}Q=P+AP-PAP=Q\quad\blacksquare
				\end{align*}
			\end{example}
	\subsection{Add}
		\begin{definition}[\texttt{[m,n] * [m,n] -> [m,n]}]
			if $A=[a_{ij}]$ and $B=[b_{ij}]$, then
			$$A+B=[a_{ij}+b_{ij}]$$
			\begin{itemize}
				\item commutativity, associativity
				\item add-identity: $0_{mn}$
				\item add-inverse: $-A$
			\end{itemize}
		\end{definition}
	\subsection{Transpose}
		\begin{definition}[\texttt{[m,n]ᵀ -> [n,m]}]
			$$A=[a_{ij}]\implies A^T=[a_{ji}]$$
		\end{definition}
		\begin{theorem}
			\begin{itemize}
				\item $(A^T)^T=A$
				\item $(kA)^T=kA^T$
				\item $(A+B)^T=A^T+B^T$
				\item $(AB)^T=B^TA^T$
				\item $\forall A\in\mathbb{R}^{1\times1}\quad A^T=A$
			\end{itemize}
		\end{theorem}
	\subsection{Scalar Mul}
		\begin{definition}[\texttt{[m,n] * k -> [m,n]}]
			if $A=[a_{ij}]$ and $k\in\mathbb{R}$, then
			$$kA=[ka_{ij}]$$
			\begin{itemize}
				\item commutativity, associativity, distributivity 
				\item mul-identity: $1$
			\end{itemize}
		\end{definition}
	\subsection{Matrix-Vector Mul}
		\begin{definition}[Matrix written as columns]
			a $m\times n$ matrix can be written as
			$$A=\begin{mat}{cccc} \vec a_1 & \vec a_2 & \cdots & \vec a_n\\ \end{mat}$$
			where $\vec a_i\in\mathbb{R}^m$
		\end{definition}
		\begin{definition}[\texttt{[m,n] * [n] -> [m]}]
			if $m\times n$ matrix $A=\begin{mat}{cccc} \vec a_1 & \vec a_2 & \cdots & \vec a_n\\ \end{mat}$ and $\vec x$ is any n-vec, then
			$$A\vec x=x_1\vec a_1+x_2\vec a_2+...+x_n\vec a_n$$
			\begin{itemize}
				\item associativity, distributivity 
			\end{itemize}
		\end{definition}
		\begin{example}[Prove $\forall \vec x\in\mathbb{R}^n\quad A\vec x=B\vec x\implies A=B$]\label{sec:standard-basis}
			for $A,B\in\mathbb{R}^{m\times n}$
			\begin{align*}
				\text{let }&A=\begin{mat}{cccc} a_1 & a_2 & \cdots & a_n \end{mat}\\
				   &B=\begin{mat}{cccc} b_1 & b_2 & \cdots & b_n \end{mat}\\
				   &\{e_i\}\text{ be standard basis of }\mathbb{R}^n
			\end{align*}
			\begin{align*}
				\text{assume }&\forall\vec x\in\mathbb{R}^n\quad A\vec x=B\vec x\\
				\therefore&\forall1\le i\le n\quad Ae_i=Be_i\implies a_i=b_i\implies A=B\quad\blacksquare
			\end{align*}
		\end{example}
		\subsubsection{System of Equations}
			\begin{theorem}[let $A\vec x=\vec{b}$ be a system of equations and rank $A=r$]
				\begin{itemize}
					\item rank $\begin{mat}{c|c}
						A & b\\
					\end{mat} = r \oplus r+1$
					\item system is \textbf{consistent} iff rank is $r$
					\item system is \textbf{inconsistent} iff rank is $r+1$
				\end{itemize}
			\end{theorem}
	\subsection{Matrix-Matrix Mul}
		\begin{definition}[\texttt{[m,n] * [n,k] -> [m,k]}]
			if $A\in\mathbb{R}^{m\times\textcolor{sky}n}$ and $B=\begin{mat}{cccc}b_1&b_2&...&b_n\end{mat}\in\mathbb{R}^{\textcolor{sky}n\times k}$, then $AB\in\mathbb{R}^{m\times k}$ and
			$$AB=A\begin{mat}{cccc}b_1&b_2&...&b_n\end{mat}=\begin{mat}{cccc}Ab_1&Ab_2&...&Ab_n\end{mat}=[c_{ij}]$$
			where $\displaystyle c_{ij}=\sum_{k=1}^na_{ik}b_{kj}$
		\end{definition}
		\begin{theorem}[Assume $a\in\mathbb{R}\quad A,B,C\in\mathbb{R}^{m\times n}$]
			\begin{itemize}
				\item $(AB)^T=B^TA^T$
				\item associativity, distributivity
					\begin{align*}
						A(BC)&=(AB)C\\
						a(AB)&=(aA)B=A(aB)\\
						ABCD&\not\equiv ACDB\\
					\end{align*}
				\item NO commutativity: $AB\not\equiv BA$\\
					$A$ \textbf{commutes} with  $B$  iff $AB=BA$
			\end{itemize}
		\end{theorem}
		\begin{example}[Prove $(AB)^T=B^TA^T$]
            let $A=[a_{ij}]$ and $B=[b_{ij}]$
            $$\therefore A^T=[a'_{ij}]\quad B^T=[b'_{ij}]\text{ where  }a'_{ij}=a_{ji}\quad b'_{ij}=b_{ji}$$
            let $B^TA^T=[c_{ij}]$
            $$\therefore c_{ij}=\sum_{k=1}^nb'_{ik}a'_{kj}=\sum_{k=1}^nb_{ki}a_{jk}=\sum_{k=1}^na_{jk}b_{ki}=(AB)_{ji}=(AB)^T_{ij}$$
		\end{example}
		\begin{example}[let $A\in\mathbb{R}^{m\times n}$ and $B\in\mathbb{R}^{n\times n}$\\Prove $A\ne0\land AB=0\implies\nexists C\in\mathbb{R}^{n\times n}$ st $BC=I$]
			assume $A\ne0$ and $AB=0$, let $C\in\mathbb{R}^{n\times n}$
			$$BC=I\implies 0=0C=(AB)C=A(BC)=AI=A$$
			which contradicts the initial assumption \blacksquare
		\end{example}
		\subsubsection{Calculate Product}
			\begin{definition}[Dot Product]
				let $a,b\in\mathbb{R}^n$
				$$a\cdot b=\sum a_ib_i$$
			\end{definition}
			\begin{algorithm}[$AB$]
				represent $A$ as rows $a_i$ and $B$ as cols $b_i$
				$$\begin{array}{c|c}
					&B\\\hline
					A&AB\\
				\end{array}=\begin{array}{c|cccc}
					&b_1&b_2&\cdots&b_n\\ \hline
					a_1&a_1\cdot b_1 & a_1\cdot b_2 & \cdots & a_1\cdot b_n\\
					a_2&a_2\cdot b_1 & a_2\cdot b_2 & \cdots & a_2\cdot b_n\\
					\vdots&\vdots & \vdots & \ddots & \vdots\\
					a_m&a_m\cdot b_1 & a_m\cdot b_2 & \cdots & a_m\cdot b_n\\
				\end{array}$$
			\end{algorithm}
			\begin{example}[Prove $\forall B\in\mathbb{R}^{2\times2}\enspace A=\begin{mat}{cc} a&0\\ 0&a\\ \end{mat}$ commutes with $B$]
				let $B=\begin{mat}{cc} b&c\\ d&e\\ \end{mat}$ for $b,c,d,e\in\mathbb{R}$
				\begin{align*}
					\begin{array}{cc|cc}
						&&b&c\\
						&&d&e\\\hline
						a&0&ab&ac\\
						0&a&ad&ae\\
					\end{array}\qquad
					\begin{array}{cc|cc}
						&&a&0\\
						&&0&a\\\hline
						b&c&ab&ac\\
						d&e&ad&ae\\
					\end{array}\\
					\therefore AB=BA=\begin{mat}{cc}
						ab&ac\\
						ad&ae\\
					\end{mat}\quad\blacksquare
				\end{align*}
			\end{example}
	\subsection{Trace}
	    \begin{definition}
			$$\mathrm{tr}(A)=\mathrm{tr}[a_{ij}]=\sum a_{ii}$$
	    \end{definition}
		\begin{theorem}
			\begin{itemize}
				\item $\mathrm{tr}(A+B)=\mathrm{tr}(A)+\mathrm{tr}(B)$ 
				\item $\mathrm{tr}(kA)=k\mathrm{tr}(A)$
				\item $\mathrm{tr}(A^T)=\mathrm{tr}(A)$
				\item $\mathrm{tr}(AB)=\mathrm{tr}(BA)$
				\item $\mathrm{tr}(AB)\not\equiv\mathrm{tr}(A)\mathrm{tr}(B)$
				\item $\mathrm{tr}(A^TB)=\mathrm{tr}(AB^T)=\sum a_{ii}b_{ii}$
			\end{itemize}
		\end{theorem}
		\begin{example}[Prove $\forall A\in\mathbb{R}^{m\times n},B\in\mathbb{R}^{n\times m}\quad\mathrm{tr}(AB)=\mathrm{tr}(BA)$]
			let $AB=C$
			\begin{align*}
				\because C_{ii}&=\sum_{j=1}^nA_{ij}B_{ji}\\\
				\therefore\mathrm{tr}(AB)=\sum_{i=1}^nC_{ii}&=\sum_{i=1}^m\sum_{j=1}^nA_{ij}B_{ji}\\
				&=\sum_{j=1}^n\sum_{i=1}^nA_{ji}B_{ij}=\mathrm{tr}(BA)
			\end{align*}
		\end{example}
	\subsection{Inverse}
		\begin{definition}
			if $A\in\mathbb{R}^{n\times n}$, then $A$ \textbf{is invertible} iff any:
			\begin{itemize}
				\item $A^T$ is invertible
				\item $\exists B$ st $AB=I$ or $\exists C$ st $CA=I$
				\item $\det A\ne0$
				\item $A=\prod E_i\equiv A\sim I$\\i.e. $A\to I$ is possible by row operations 
				\item $\ker(A)=\{0\}$\\i.e. $A\vec x=0$ has only the solution $\vec x=0$ 
				\item $\forall\vec b\quad A\vec x=\vec b$ has solutions 
				\item $\mathrm{rank}(A)=n$
			\end{itemize}
			if $AB=I=CA$ for all square matrices, then
			$$B=C=A^{-1}$$
		\end{definition}
		\begin{theorem}
			[Assume $A$ is invertible/has an inverse]
			\begin{itemize}
				\item $(A^T)^{-1}=(A^{-1})^T$ 
				\item $(A_1A_2\cdots A_n)^{-1}=A_n^{-1}\cdots A_2^{-1}A_1^{-1}$
				\item $(A^k)^{-1}=(A^{-1})^k$
				\item $(aA)^{-1}=\sfrac1aA^{-1}$ for $a\ne0$
			\end{itemize}
		\end{theorem}
		\begin{algorithm}[Find Inverse]
			use Gaussian Algorithm to
			$$\begin{mat}{c|c} A & I\\ \end{mat}\rightarrow\begin{mat}{c|c} I & A^{-1}\\ \end{mat}$$
		\end{algorithm}
		\begin{example}
			[Prove if $A,B$ are invertible, then $(AB)^{-1}=B^{-1}A^{-1}$]
			if $(L)^{-1}=R$, show $LR=I=RL$
			\begin{align*}
				(AB)(B^{-1}A^{-1})=A(BB^{-1})A^{-1}=AIA^{-1}=AA^{-1}=I\\
				(B^{-1}A^{-1})(AB)=B^{-1}(A^{-1}A)B=B^{-1}IB=B^{-1}B=I\quad\blacksquare\\
			\end{align*}
		\end{example}
		\begin{example}
			[Prove $AB$ is invertible $\implies A$ and $B$ are both invertible] 
			assume $AB\in\mathbb{R}^{n\times n}$ is invertible
			\begin{align*}
				\therefore\exists C\in\mathbb{R}^{n\times n}\text{ st }(AB)C&=C(AB)=I\\
				A(BC)&=I\implies A\text{ is invertible }\\
				(CA)B&=I\implies B\text{ is invertible }\\
			\end{align*}
		\end{example}
		\subsubsection{Solution of a System}
			\begin{theorem}[Find Unique Solution of a System]
				if a system of $n$ eqns and  $n$ vars is written as
				$$A\vec x=\vec b$$
				then 
				$$A\text{ is invertible }\implies \vec x=A^{-1}\vec b$$
			\end{theorem}
	\subsection{Elementary Matrices}
		\begin{definition}
			an $n\times n$ matrix $E$ is \textbf{elementary} if $E$ can be produced from a single elementary row operation on $I$
		\end{definition}
		\begin{theorem}[Lemma]
			suppose an elementary operation \texttt{op()} and $A\in\mathbb{R}^{m\times n}$
			\begin{itemize}
				\item $\text{op}(A)=\text{op}(I_m)A=EA$
				\item $\text{op}(E^{-1})=I$ where $E^{-1}$ is also elementary and is the reverse operation
			\end{itemize}
		\end{theorem}
		\subsubsection{Theorem}
			\begin{theorem}
				[Suppose $A\in\mathbb{R}^{m\times n}$ and $A\to R$ by elem row ops]
				\begin{itemize}
					\item $A\to R\equiv\begin{mat}{cc} A&I_m\\ \end{mat}\to\begin{mat}{cc} R&U\\ \end{mat}$
					\item $R=UA$ for $U\in\mathbb{R}^{m\times m}$ is invertible
					\item $U=E_n\cdots E_2E_1$ where $E_i$ is the $i$th row operation
					\item $B\in\mathbb{R}^{n\times n}$ is invertible iff $B=\prod E_i$
					\item $A=UR\implies A\overset r\sim R$ \textbf{row equivalent}
				\end{itemize}
			\end{theorem} 
			\begin{example}[Prove $A\overset r\sim B\Rightarrow B\overset r\sim A$ and $A\overset r\sim B\overset r\sim C\Rightarrow A\overset r\sim C$]
				\begin{enumerate}
					\item assume $A\overset r\sim B$
						\begin{align*}
							\therefore&A=UB\text{ for invertible }U\\
							\therefore&B=U^{-1}A\implies B\overset r\sim A\quad\blacksquare\\
						\end{align*}
					\item assume $A\overset r\sim B\overset r\sim C$
						\begin{align*}
							\therefore&A=UB\quad B=VC\text{ for invertible }U,V\\
							\therefore&A=UVC=(UV)C\implies A\overset r\sim C\quad\blacksquare
						\end{align*}
				\end{enumerate}
			\end{example}
			\begin{example}
				[let $A\in\mathbb{R}^{m\times n}$ and $B\in\mathbb{R}^{n\times m}$\\Prove $m>n\implies AB$ is not invertible]
				\begin{align*}
					m>n\implies&B\vec x=0\text{ has nontrivial solution }\\
					\therefore&\,\exists\vec x\ne0\text{ st }B\vec x=0\\
					\therefore&\,A(B\vec x)=AB\vec x=0\text{ has nontrivial solution }\\
					\therefore&\,AB\text{ is not invertible }\quad\blacksquare
				\end{align*}
			\end{example}
	\subsection{Linear Map}
	    \begin{definition}[$\forall x,y\in\mathbb{R}^n$]
	    	\begin{itemize}
				\item if $T$ is linear, then $T(x)=T_A(x)=Ax$ where
					$$A=\begin{mat}{cccc}
						T(e_1)&T(e_2)&\cdots&T(e_n)\\
					\end{mat}$$
				\item $T(x+y)=T(x)+T(y)$
				\item $T(ax)=aT(x)$
				\item $T=A\land S=B\implies(T\circ S)(x)=(AB)x$
	    	\end{itemize}
	    \end{definition}
		\begin{theorem}[Linearity]
			if $T:\mathbb{R}^n\to\mathbb{R}^m$ is a linear transformation, then
			$$T\left(\sum a_i\vec{x_i}\right)=\sum a_iT(\vec{x_i})$$
			proven using induction
		\end{theorem}
		\begin{definition}[Kernel and Image for $T:\mathbb{R}^n\to\mathbb{R}^m$]
			\begin{align*}
				\ker(T)&=\{x\in\mathbb{R}^n:T(x)=0\}\\
				\mathrm{im}(T)&=\{T(x):x\in\mathbb{R}^n\}
			\end{align*}
			\begin{itemize}
				\item $x_1,x_2\in\ker(T)\implies ax_1+bx_2\in\ker(T)$ 
				\item $x_1,x_2\in\mathrm{im}(T)\implies ax_1+bx_2\in\mathrm{im}(T)$ 
			\end{itemize}
		\end{definition}
		\begin{theorem}[Geometrical Linear Transformatins]
			they are all $T:\mathbb{R}^2\to\mathbb{R}^2$
			\begin{align*}
				&R_\theta=\begin{mat}{cc}
					\cos\theta&-\sin\theta\\
					\sin\theta&\cos\theta\\
				\end{mat}\tag*{(rotate anticlock-wise)}\\
				&Q_m=\frac1{1+m^2}\begin{mat}{cc}
					1-m^2&2m\\
					2m&m^2-1\\
				\end{mat}\tag*{(reflect on $y=mx$)}
			\end{align*}
		\end{theorem}
	\subsection{Practice}
		\begin{example}[let $A\in\mathbb{R}^{n\times n}$ st $\vec x\ne0\implies\vec x^TA\vec x>0$\\
			Show $\forall $skew-symmetric $B\in\mathbb{R}^{n\times n},\vec b\in\mathbb{R}^n$
			$$(A+B)\vec x=\vec b$$ as a unique solution]
			\begin{align*}
				\because\,&(A+B)\vec x=\vec b\text{ has a unique solution iff }(A+B)\vec x=0\implies\vec x=0\\
				&(\vec x\ne0\implies\vec x^TA\vec x>0)\equiv(\vec x^TA\vec x\le0\implies\vec x=0)\\
				\therefore\,&\text{given }\vec x^TA\vec x\le0\implies\vec x=0\text{ and }B^T=-B\\
				&\text{show }(A+B)\vec x=0\implies\vec x=0\\
			\end{align*}
			\begin{enumerate}
				\item show $\vec x^TB\vec x=0$
				\begin{align*}
					\vec x^TB\vec x&=(\vec x^TB\vec x)^T\tag*{($A^T=A\quad\forall A\in\mathbb{R}^{1\times1}$)}\\
					\vec x^TB^T\vec x&=-\vec x^TB\vec x\\
					\iff\vec x^TB\vec x&=-\vec x^TB\vec x\\
					2\vec x^TB\vec x&=0\\
					\vec x^TB\vec x&=0\\
				\end{align*}
				\item show $(A+B)\vec x=0\implies\vec x=0$
				\begin{align*}
					(A+B)\vec x=A\vec x+B\vec x&=0\tag*{(distributivity)}\\
					\vec x^TA\vec x+\vec x^TB\vec x&=0\tag*{(left mul by $\vec x^T$)}\\
					\vec x^TA\vec x&=0\tag*{(from 1)}\\
					\vec x^TA\vec x&\le0\implies\vec x=0\tag*{(given)}
				\end{align*}
			\end{enumerate}
		\end{example}
\section{Eigenvectors}
    \subsection{Cofactor Expansion}
	    \begin{definition}[$\det(A)$ for $A\in\mathbb{R}^{2\times2}$ and $B\in\mathbb{R}^{3\times3}$]
			\begin{align*}
				\det(A)&=\begin{vmat}{cc}a&b\\c&d\end{vmat}=ad-bc\\
				\det(B)&=\begin{vmat}{ccc} a&b&c\\ d&e&f\\ g&h&i\\ \end{vmat}=a\begin{vmat}{cc} e&f\\ h&i\\ \end{vmat}-b\begin{vmat}{cc} d&f\\ g&i\\ \end{vmat}+c\begin{vmat}{cc} d&e\\ g&h\\ \end{vmat}\\
			\end{align*}
	    \end{definition}
		\begin{definition}[Cofactor expansion of $\det(A)$ along row $1$]
			for $A=[a_{ij}]\in\mathbb{R}^{n\times n}$
			\begin{align*}
				\det(A)=|A|&=\sum a_{1i}c_{1i}(A)\\
				c_{ij}(A)&=(-1)^{i+j}|A_{ij}|\\
				A_{ij}\text{ is the matrix of }&A\text{ after removing row }i\text{ and col }j\\
			\end{align*}
		\end{definition}
		\begin{theorem}[let $A\in\mathbb{R}^{n\times n}$ and $u\in\mathbb{R}$]
			\begin{itemize}
				\item if $\exists$ a row/col is a \textbf{linear combination} of any other rows/cols\\
					then $\color{\colo}|A|=0$ 
				\item if $B$ is $A$ but 2 rows/cols are swapped\\
					then $\color{\colo}|B|=-|A|$
				\item if $B$ is $A$ but a row/col\texttt{ *= c}\\
					then $\color{\colo}|B|=c|A|$
				\item \texttt{Aᵢ += c} $\implies$ $c|A|$
				\item \texttt{Aᵢ += u * Aⱼ} for \texttt{i != j} $\implies$ \textbf{keeps} $|A|$
			\end{itemize}
		\end{theorem}

	\subsection{Determinant and Matrix Inverse}
	    \subsubsection{Determinant}
		    \begin{theorem}
				\begin{itemize}
					\item $|I|=1$
				    \item $|cA|=c^n|A|$
				    \item $|A^T|=|A|$ for $A\in\mathbb{R}^{n\times n}$
				    \item $|AB|=|A||B|$
				    \item $|A^{-1}|=\frac1{|A|}$ for invertible $A$
				\end{itemize}
		    \end{theorem}
			\begin{example}[find $|A|$ if $A^3=A$]
				\begin{align*}
					A^3&=A\implies|A^3|=|A|\implies|A|^3-|A|=0\\
					|A|(|A|^2-1)&=0\\
					\therefore |A|&=0,\pm1
				\end{align*}
			\end{example}
		\subsubsection{Adjugate}
		    \begin{definition}
				\begin{align*}
					\mathrm{adj}(A)&=[c_{ij}(A)]^T\\
					A=\begin{mat}{cc} a&b\\ c&d\\ \end{mat}&\implies\mathrm{adj}(A)=\begin{mat}{cc} d&-b\\ -c&a\\ \end{mat}
				\end{align*}
		    \end{definition}
			\begin{theorem}
				\begin{itemize}
					\item $A(\mathrm{adj}\,A)=|A|I=(\mathrm{adj}\,A)A$
					\item $A^{-1}=\frac1{|A|}\mathrm{adj}\,A$
				\end{itemize}
			\end{theorem}
	\subsection{Eigenvalue}
	    \begin{definition}[$\lambda$]
			$\lambda$ is \textbf{an eigenvalue} of $A\in\mathbb{R}^{n\times n}$ if for $\vec x\ne0\in\mathbb{R}^n$
			$$A\vec x=\lambda\vec x\equiv|\lambda I-A|=0$$
			and $\vec x$ is an $\lambda$-eigenvector
	    \end{definition}
		\begin{definition}[$c_A(\lambda)$]
			the \textbf{characteristic polynomial} of $A$ is 
			$$c_A(\lambda)=|\lambda I-A|=(\lambda-\lambda_1)(\lambda-\lambda_2)\cdots(\lambda-\lambda_n)$$
			\begin{itemize}
				\item \textbf{eigenvalues} of $A$ is $\lambda$ when
					$$c_A(\lambda)=0$$
				\item $\lambda$\textbf{-eignvectors} is $\vec x\ne0$ when
					$$(\lambda I-A)\vec x=0$$
			\end{itemize}
		\end{definition}
		\subsubsection{A-Invariance}
		    \begin{definition}[for $A\in\mathbb{R}^{n\times n}$ and $\vec x\ne0\in\mathbb{R}^n$]
				\begin{itemize}
					\item $L_x=\{\lambda\vec x:\lambda\in\mathbb{R}\}$ is \textbf{$A$-invariant} if $\vec x\in L\implies A\vec x\in L$
					\item $\vec x$ is an eigenvec of $A$ iff $L_x$ is $A$-invariant
				\end{itemize}
		    \end{definition}
		\subsubsection{Multiplicity}
		    \begin{definition}
				an eigenval $\lambda$ has \textbf{multiplicity} $n$ if it occurs $n$ items as root of $c_A(x)$
				\begin{itemize}
					\item $(\lambda-1)^2\implies\lambda_1=1$ with multiplicity $2$
				\end{itemize}
		    \end{definition}
		\subsubsection{Diagonalisation}
		    \begin{definition}[Power of Matrix]
		        if $A=PDP^{-1}$ then $A^k=PD^kP^{-1}$ for $k\in\mathbb{Z}^+$
		    \end{definition}
		    \begin{definition}[Diagonal Matrix]
		        $$D=\mathrm{diag}(\lambda_1,\lambda_2,\cdots,\lambda_n)=\begin{mat}{cccc} \lambda_1e_1&\lambda_2e_2&\cdots&\lambda_ne_n\\ \end{mat}$$
		    \end{definition}
			\begin{definition}
				$A\in\mathbb{R}^{n\times n}$ is \textbf{diagonalisable} if
				$$\exists P\in\mathbb{R}^{n\times n}\text{ st }P^{-1}AP=D$$
				where the invertible $P$ is a \textbf{diagonalising matrix} of $A$
			\end{definition}
			\begin{theorem}
			    for $A\in\mathbb{R}^{n\times n}$
				\begin{itemize}
					\item $A$ is diagable iff $A$ has eigenvecs $\vec x_i$ st $P=\begin{mat}{c}\vec x_i\end{mat}$ is invertible
					\item if diagable, $P^{-1}AP=\mathrm{diag}(\lambda_1,\cdots,\lambda_n)$ where $A\vec x_i=\lambda_i\vec x_i$
				\end{itemize}
			\end{theorem}
			\begin{algorithm}[Diagonise $A\in\mathbb{R}^{n\times n}$]
				\begin{enumerate}
					\item find distinct eigenvals $\lambda$ of $A$
					\item for each $\lambda$, find basic eigenvec $\vec x$ from $(\lambda I-A)\vec x=0$
					\item $A$ is diagable iff $\exists n$ eigenvecs
					\item if diagable, then $P=\begin{mat}{c}\vec x_i\end{mat}$ st $P^{-1}AP$ is diagonal
				\end{enumerate}
			\end{algorithm}
\end{document}
